% Options for packages loaded elsewhere
\PassOptionsToPackage{unicode}{hyperref}
\PassOptionsToPackage{hyphens}{url}
\PassOptionsToPackage{dvipsnames,svgnames,x11names}{xcolor}
\documentclass[
  10pt,
  fontset=fandol]{ctexart}
\usepackage{xcolor}
\usepackage[margin=16mm]{geometry}
\usepackage{amsmath,amssymb}
\setcounter{secnumdepth}{-\maxdimen} % remove section numbering
\usepackage{iftex}
\ifPDFTeX
  \usepackage[T1]{fontenc}
  \usepackage[utf8]{inputenc}
  \usepackage{textcomp} % provide euro and other symbols
\else % if luatex or xetex
  \usepackage{unicode-math} % this also loads fontspec
  \defaultfontfeatures{Scale=MatchLowercase}
  \defaultfontfeatures[\rmfamily]{Ligatures=TeX,Scale=1}
\fi
\usepackage{lmodern}
\ifPDFTeX\else
  % xetex/luatex font selection
\fi
% Use upquote if available, for straight quotes in verbatim environments
\IfFileExists{upquote.sty}{\usepackage{upquote}}{}
\IfFileExists{microtype.sty}{% use microtype if available
  \usepackage[]{microtype}
  \UseMicrotypeSet[protrusion]{basicmath} % disable protrusion for tt fonts
}{}
\makeatletter
\@ifundefined{KOMAClassName}{% if non-KOMA class
  \IfFileExists{parskip.sty}{%
    \usepackage{parskip}
  }{% else
    \setlength{\parindent}{0pt}
    \setlength{\parskip}{6pt plus 2pt minus 1pt}}
}{% if KOMA class
  \KOMAoptions{parskip=half}}
\makeatother
\usepackage{longtable,booktabs,array}
\newcounter{none} % for unnumbered tables
\usepackage{calc} % for calculating minipage widths
% Correct order of tables after \paragraph or \subparagraph
\usepackage{etoolbox}
\makeatletter
\patchcmd\longtable{\par}{\if@noskipsec\mbox{}\fi\par}{}{}
\makeatother
% Allow footnotes in longtable head/foot
\IfFileExists{footnotehyper.sty}{\usepackage{footnotehyper}}{\usepackage{footnote}}
\makesavenoteenv{longtable}
\setlength{\emergencystretch}{3em} % prevent overfull lines
\providecommand{\tightlist}{%
  \setlength{\itemsep}{0pt}\setlength{\parskip}{0pt}}
\usepackage{amsmath,amssymb,mathtools,needspace}
\xeCJKDeclareCharClass{CJK}{"2032,"0400->"04FF,"2160->"216B,"2460->"2473,"25A0,"25C7,"25CB,"25CF}
\IfFontExistsTF{Arial Unicode MS}{\xeCJKsetup{AutoFallBack=true}\setCJKfallbackfamilyfont{\CJKrmdefault}{Arial Unicode MS}}{}
\setcounter{secnumdepth}{0}
\setlength{\emergencystretch}{2em}
\allowdisplaybreaks
\usepackage{bookmark}
\IfFileExists{xurl.sty}{\usepackage{xurl}}{} % add URL line breaks if available
\urlstyle{same}
\hypersetup{
  pdftitle={Newton 法应用举例},
  colorlinks=true,
  linkcolor={Maroon},
  filecolor={Maroon},
  citecolor={Blue},
  urlcolor={Blue},
  pdfcreator={LaTeX via pandoc}}

\title{Newton 法应用举例}
\author{}
\date{}

\begin{document}
\maketitle

\subsubsection{6.3.4 Newton 法应用举例}\label{newton-ux6cd5ux5e94ux7528ux4e3eux4f8b}

对于给定正数 \(a\)，应用 Newton 法解二次方程 \(x^2-a=0\)，可导出求开方值 \(\sqrt{a}\) 的计算程序

\[
x_{k+1}=\frac{1}{2}\left(x_k+\frac{a}{x_k}\right).
\tag{6.3.7}
\]

下面证明这种迭代公式对于任意初值 \(x_0>0\) 都是收敛的。

事实上，对式（6.3.7）施行配方手续，易知

\[
x_{k+1}-\sqrt{a}=\frac{1}{2x_k}(x_k-\sqrt{a})^2,\quad
x_{k+1}+\sqrt{a}=\frac{1}{2x_k}(x_k+\sqrt{a})^2.
\]

将以上两式相除得

\begin{quote}
校注（agent 补充）：原页步 4 的替换三元组印为 \((x_1,f_1,f')\)，第三项未印下标 \(1\)，转录保留原式。结合步 2 的定义及被替换的 \(f'_0\)，这里疑应为 \(f'_1\)。这是原扫描上的疑似漏印，非 OCR 改写；局部原图见 \href{../assets/ch06/detail-p165-step4.png}{步 4 裁切}。
\end{quote}

\begin{quote}
校注（agent 补充）：本页首段``单根''推出``平方收敛''的原文亦予保留。若按前页定义 6.2 要求精确收敛阶的极限常数非零，则还需 \(f''(x^*)\ne0\) 才能断言恰为二阶；在通常光滑性条件下，单根保证的是局部至少二阶收敛。例如 \(f(x)=x+x^3\) 在单根 \(0\) 附近的 Newton 迭代为 \(x_{k+1}=2x_k^3/(1+3x_k^2)\)，实际为三阶。这是原文收敛阶表述的适用条件说明。
\end{quote}

\href{../../source-images/pdf-166.jpeg}{原页}

\[
\frac{x_{k+1}-\sqrt{a}}{x_{k+1}+\sqrt{a}}=\left[\frac{x_k-\sqrt{a}}{x_k+\sqrt{a}}\right]^2.
\]

据此反复递推，有

\[
\frac{x_k-\sqrt{a}}{x_k+\sqrt{a}}=\left[\frac{x_0-\sqrt{a}}{x_0+\sqrt{a}}\right]^{2^k},
\tag{6.3.8}
\]

记 \(q=\dfrac{x_0-\sqrt{a}}{x_0+\sqrt{a}}\)，整理式（6.3.8），得

\[
x_k-\sqrt{a}=2\sqrt{a}\frac{q^{2^k}}{1-q^{2^k}}.
\]

对于任意 \(x_0>0\)，总有 \(|q|<1\)，故由上式推知，当 \(k\to\infty\) 时 \(x_k\to\sqrt{a}\)，即迭代过程恒收敛。

\textbf{例 6.8} 求 \(\sqrt{115}\)。

\textbf{解} 取初值 \(x_0=10\)，对 \(a=115\) 按式（6.3.7）迭代 3 次便得到精度为 \(10^{-6}\) 的结果（见表 6.8）。

表 6.8

{\def\LTcaptype{none} % do not increment counter
\begin{longtable}[]{@{}ll@{}}
\toprule\noalign{}
\(k\) & \(x_k\) \\
\midrule\noalign{}
\endhead
\bottomrule\noalign{}
\endlastfoot
\(0\) & \(10\) \\
\(1\) & \(10.750\,000\) \\
\(2\) & \(10.723\,837\) \\
\(3\) & \(10.723\,805\) \\
\(4\) & \(10.723\,805\) \\
\end{longtable}
}

由于式（6.3.7）对于任意初值 \(x_0>0\) 均收敛，并且收敛的速度很快，因此，可取确定的初值如 \(x_0=1\) 编制通用程序。用这个通用程序求 \(\sqrt{115}\)，也只要迭代 7 次便得到了上面的结果 \(10.723\,805\)。

再举一个例子。

对于给定的正数 \(a\)，应用 Newton 法于方程 \(1/x-a=0\)，可导出求 \(1/a\) 而不用除法的计算程序，即

\[
x_{k+1}=x_k(2-ax_k).
\]

这个算法有实际意义，早期设计电子计算机时，为节省硬件设备，曾运用这种技术避开除法操作的设置。

下面证明这一算法当初值 \(x_0\) 满足 \(0<x_0<2/a\) 时是收敛的。事实上，由于

\[
x_{k+1}-\frac{1}{a}=x_k(2-ax_k)-\frac{1}{a}=-a\left(x_k-\frac{1}{a}\right)^2,
\]

因而，对 \(r_k=1-ax_k\) 有递推公式 \(r_{k+1}=r_k^2\)。据此反复递推，有

\[
r_k=r_0^{2^k}.
\]

若初值满足 \(0<x_0<\dfrac{2}{a}\)，则对 \(r_0=1-ax_0\) 有 \(|r_0|<1\)。这时将有 \(r_k\to0\)，从而迭代收敛。

\subsection{本节覆盖原页的校注记录}\label{ux672cux8282ux8986ux76d6ux539fux9875ux7684ux6821ux6ce8ux8bb0ux5f55}

下列记录按原页关联，含同页相邻小节的校注；计算时同时读取上文校注和完整原页。

\begin{itemize}
\tightlist
\item
  \texttt{ch06-E002}：原书 PDF 页 165
\item
  \texttt{ch06-E008}：原书 PDF 页 165
\end{itemize}

\end{document}
